\textbf{Phần 1. Trắc nghiệm lựa chọn}
\setcounterpageref{ex}{0}

\begin{ex} %[2H5V28]
Một đường gấp khúc gồm $3$ đoạn thẳng có độ dài lần lượt là $25\text{ cm}, 38\text{ cm}, 27\text{ cm}$. Chiều dài đường gấp khúc đó là
\choice
{$88\text{ cm}$}
{$89\text{ cm}$}
{$91\text{ cm}$}
{\True $90\text{ cm}$}
\loigiai{
$25 + 38 + 27 = 90\text{ (cm)}$.
\textbf{Đáp án của câu hỏi là D}
}
\end{ex}

\begin{ex} %[2H5V28]
Một hình tam giác có chu vi $48\text{ cm}$, biết tổng độ dài hai cạnh là $31\text{ cm}$. Độ dài cạnh còn lại là
\choice
{$16\text{ cm}$}
{\True $17\text{ cm}$}
{$18\text{ cm}$}
{$15\text{ cm}$}
\loigiai{
$48 - 31 = 17\text{ (cm)}$.
\textbf{Đáp án của câu hỏi là B}
}
\end{ex}

\begin{ex} %[2H5V28]
Trong các hình dưới đây, hình nào có số cạnh nhiều nhất?
\choice
{Hình tam giác}
{Đoạn thẳng}
{\True Hình tứ giác}
{Đường gấp khúc có $2$ đoạn thẳng}
\loigiai{
Hình tứ giác có 4 cạnh, nhiều nhất trong các phương án.
\textbf{Đáp án của câu hỏi là C}
}
\end{ex}

\begin{ex} %[2H5V28]
Đoạn thẳng $AB$ dài $42\text{ cm}$. Điểm $M$ nằm giữa $A$ và $B$ sao cho $AM = 18\text{ cm}$. Đoạn thẳng $MB$ dài bao nhiêu?
\choice
{\True $24\text{ cm}$}
{$22\text{ cm}$}
{$23\text{ cm}$}
{$25\text{ cm}$}
\loigiai{
$MB = 42 - 18 = 24\text{ (cm)}$.
\textbf{Đáp án của câu hỏi là A}
}
\end{ex}

\begin{ex} %[2H5V28]
Cho hình vuông có cạnh $15\text{ cm}$. Chu vi của hình vuông đó là
\choice
{$50\text{ cm}$}
{$55\text{ cm}$}
{$65\text{ cm}$}
{\True $60\text{ cm}$}
\loigiai{
Chu vi: $15 + 15 + 15 + 15 = 60\text{ (cm)}$.
\textbf{Đáp án của câu hỏi là D}
}
\end{ex}

\begin{ex} %[2H5V28]
Một con kiến bò quanh mép một cái bánh hình tam giác có các cạnh lần lượt là $8\text{ cm}, 12\text{ cm}, 15\text{ cm}$. Quãng đường kiến bò được là
\choice
{$33\text{ cm}$}
{\True $35\text{ cm}$}
{$34\text{ cm}$}
{$36\text{ cm}$}
\loigiai{
Quãng đường: $8 + 12 + 15 = 35\text{ (cm)}$.
\textbf{Đáp án của câu hỏi là B}
}
\end{ex}

\begin{ex} %[2H5V28]
Một đoạn dây đồng dài $72\text{ cm}$ được uốn thành một hình tứ giác có $3$ cạnh lần lượt dài $16\text{ cm}, 19\text{ cm}, 21\text{ cm}$. Cạnh thứ tư dài bao nhiêu?
\choice
{$15\text{ cm}$}
{$17\text{ cm}$}
{\True $16\text{ cm}$}
{$18\text{ cm}$}
\loigiai{
Cạnh thứ tư: $72 - (16 + 19 + 21) = 72 - 56 = 16\text{ (cm)}$.
\textbf{Đáp án của câu hỏi là C}
}
\end{ex}

\begin{ex} %[2H5V28]
Hình gồm $5$ điểm phân biệt thẳng hàng có bao nhiêu đoạn thẳng?
\choice
{\True $10$ đoạn thẳng}
{$8$ đoạn thẳng}
{$12$ đoạn thẳng}
{$6$ đoạn thẳng}
\loigiai{
$5 \times 4 : 2 = 10$ đoạn thẳng.
\textbf{Đáp án của câu hỏi là A}
}
\end{ex}

\begin{ex} %[2H5V28]
Điền số thích hợp: Đường gấp khúc dài $1\text{ m}$ gồm $2$ đoạn thẳng. Đoạn thứ nhất dài $65\text{ cm}$, đoạn thứ hai dài
\choice
{$45\text{ cm}$}
{$25\text{ cm}$}
{$30\text{ cm}$}
{\True $35\text{ cm}$}
\loigiai{
Đổi $1\text{ m} = 100\text{ cm}$. Đoạn thứ hai: $100 - 65 = 35\text{ (cm)}$.
\textbf{Đáp án của câu hỏi là D}
}
\end{ex}

\begin{ex} %[2H5V28]
Hình tứ giác $ABCD$ có chu vi $84\text{ cm}$. Nếu tăng độ dài mỗi cạnh thêm $2\text{ cm}$ thì chu vi của hình tứ giác mới là
\choice
{$86\text{ cm}$}
{\True $92\text{ cm}$}
{$88\text{ cm}$}
{$90\text{ cm}$}
\loigiai{
Có 4 cạnh, mỗi cạnh tăng 2 cm nên chu vi tăng: $2 \times 4 = 8\text{ cm}$.\nChu vi mới: $84 + 8 = 92\text{ (cm)}$.
\textbf{Đáp án của câu hỏi là B}
}
\end{ex}

\begin{ex} %[2H5V28]
Một cái khung tranh hình chữ nhật có chiều dài $30\text{ cm}$, chiều rộng $20\text{ cm}$. Chu vi khung tranh đó là
\choice
{$80\text{ cm}$}
{$90\text{ cm}$}
{\True $100\text{ cm}$}
{$50\text{ cm}$}
\loigiai{
Chu vi: $30 + 20 + 30 + 20 = 100\text{ (cm)}$.
\textbf{Đáp án của câu hỏi là C}
}
\end{ex}

\begin{ex} %[2H5V28]
Hình tam giác có độ dài các cạnh $a, b, c$. Biết $a = 15\text{ cm}, b = 20\text{ cm}$, chu vi tam giác bằng $50\text{ cm}$. Độ dài cạnh $c$ là
\choice
{\True $15\text{ cm}$}
{$16\text{ cm}$}
{$14\text{ cm}$}
{$18\text{ cm}$}
\loigiai{
$c = 50 - (15 + 20) = 15\text{ (cm)}$.
\textbf{Đáp án của câu hỏi là A}
}
\end{ex}

\textbf{Phần 2. Trắc nghiệm đúng sai}
\setcounterpageref{ex}{0}

\begin{ex} %[2H5V28]
Xét tính đúng sai của các bài toán hình học ôn tập Chủ đề 5:
\choiceTF
{\True Hình chữ nhật có 4 góc vuông và 4 cạnh.}
{Một đường gấp khúc có 3 đoạn thẳng thì có đúng 3 điểm.}
{\True Độ dài đường gấp khúc bằng tổng độ dài các đoạn thẳng thành phần.}
{\True Một hình vuông có chu vi 40 cm thì cạnh của nó dài 10 cm.}
\loigiai{
\begin{itemize}
    \item Ý a) Đúng: Đúng theo đặc điểm hình chữ nhật.
    \item Ý b) Sai: Đường gấp khúc 3 đoạn thẳng có 4 điểm đầu mút.
    \item Ý c) Đúng: Đúng.
    \item Ý d) Đúng: $40 : 4 = 10\text{ cm}$. Đúng.
\end{itemize}
\begin{itemize}
    \item a) \textbf{Đúng}.
    \item b) \textbf{Sai}.
    \item c) \textbf{Đúng}.
    \item d) \textbf{Đúng}.
\end{itemize}
}
\end{ex}

\begin{ex} %[2H5V28]
Cho đường gấp khúc ABCD có AB = 15 cm, BC = 25 cm, CD = 35 cm:
\choiceTF
{\True Chiều dài đường gấp khúc ABCD là 75 cm.}
{\True Đoạn CD dài hơn đoạn AB là 20 cm.}
{Đoạn BC dài gấp đôi đoạn AB.}
{Nối điểm A với điểm D thì ta được một hình tam giác.}
\loigiai{
\begin{itemize}
    \item Ý a) Đúng: $15 + 25 + 35 = 75\text{ cm}$.
    \item Ý b) Đúng: $35 - 15 = 20\text{ cm}$.
    \item Ý c) Sai: 25 không gấp đôi 15. Sai.
    \item Ý d) Sai: Nối A với D tạo thành hình tứ giác ABCD.
\end{itemize}
\begin{itemize}
    \item a) \textbf{Đúng}.
    \item b) \textbf{Đúng}.
    \item c) \textbf{Sai}.
    \item d) \textbf{Sai}.
\end{itemize}
}
\end{ex}

\textbf{Phần 3. Trả lời ngắn}
\setcounterpageref{ex}{0}

\begin{ex} %[2H5V28]
Tính chu vi của hình tam giác có độ dài ba cạnh lần lượt là 19 cm, 28 cm, 33 cm.
\shortans{80}
\loigiai{
\textbf{Đáp số:} $80$\\
$19 + 28 + 33 = 80\text{ cm}$.
}
\end{ex}

\begin{ex} %[2H5V28]
Một hình tứ giác có chu vi 90 cm. Tổng độ dài ba cạnh đầu là 68 cm. Tính độ dài cạnh thứ tư.
\shortans{22}
\loigiai{
\textbf{Đáp số:} $22$\\
$90 - 68 = 22\text{ cm}$.
}
\end{ex}

\begin{ex} %[2H5V28]
Đường gấp khúc MNPQ dài 64 cm. Biết MN = 18 cm, NP = 24 cm. Tính độ dài đoạn thẳng PQ.
\shortans{22}
\loigiai{
\textbf{Đáp số:} $22$\\
$PQ = 64 - (18 + 24) = 64 - 42 = 22\text{ cm}$.
}
\end{ex}

\begin{ex} %[2H5V28]
Cho hình vuông có chu vi bằng 36 cm. Tính độ dài một cạnh của hình vuông đó.
\shortans{9}
\loigiai{
\textbf{Đáp số:} $9$\\
$36 : 4 = 9\text{ cm}$.
}
\end{ex}

\begin{ex} %[2H5V28]
Một đoạn dây thép dài 85 cm, người ta cắt đi 38 cm để uốn khung tranh. Hỏi đoạn dây thép còn lại dài bao nhiêu xăng-ti-mét?
\shortans{47}
\loigiai{
\textbf{Đáp số:} $47$\\
$85 - 38 = 47\text{ cm}$.
}
\end{ex}

\begin{ex} %[2H5V28]
Một lối đi quanh vườn hình chữ nhật có chiều dài 25 m, chiều rộng 15 m. Tính chu vi lối đi quanh vườn đó.
\shortans{80}
\loigiai{
\textbf{Đáp số:} $80$\\
Chu vi: $25 + 15 + 25 + 15 = 80\text{ m}$.
}
\end{ex}
